\section{Threats to Validity}

The controlled campaign uses generated codebases and deterministic patches. The ground truth is injected by construction and does not represent natural project diversity.

The external campaign improves external validity by using five real open-source codebases and a first-pass static adjudication of 82 findings. This adjudication verifies static codebase evidence but is not an independent audit of every CVE status and does not validate application-level exploitability. Recall over naturally occurring external vulnerabilities is not claimed because no complete independent ground truth exists for the external codebases.

The injected external campaign provides complete ground truth only for the injected findings, not for every vulnerability that may naturally exist in the base codebases. External campaigns use an explicit build-skip mode: they measure security-pipeline behavior, but they do not prove that remediated third-party applications still pass their full project-specific test suites. The NodeGoat npm remediation uses non-breaking \texttt{npm audit fix}; remaining findings require breaking upgrades, replacement of unmaintained packages, or manual redesign.

The toolchain is time-sensitive because scanner databases, npm advisory metadata, and package-manager resolution behavior change. The reported run fixes tool versions and integrity metadata, but future runs may produce different vulnerability counts as advisory databases evolve.

The controlled precision estimate treats scanner findings not linked to injected ground truth as unknown or false-positive candidates. The external adjudication improves this point for the reviewed findings, but it remains a static first-pass review rather than a blinded multi-reviewer audit. Stronger claims about external false-positive rates require independent reviewers and a broader corpus.

Remote CI execution and real pull-request creation were performed for the reported validation. The remote validation covers compilation, unit tests, OPA policy tests, workflow syntax validation, automated remediation-branch creation, user-pushed remediation review flow, passing pull-request CI, owner merge acceptance, measured merge latency, and passing post-merge CI. This remains a functional validation rather than a measurement of independent human review quality, branch protection policy, multi-reviewer governance, or production deployment.
